\documentclass[12pt]{article}
\usepackage{amsmath, amssymb, amsthm}
\usepackage[margin=1in]{geometry}
\title{Higher Math Elective HW1}
\author{Anatolii Larkin}

\begin{document}

\maketitle

\begin{enumerate}

    \item Using the truth tables we’ve discussed in class, evaluate the following expressions in classical Boolean logic with two truth values \(B = \{T, F\}\):

    \begin{enumerate}
        \item[(a)] \((T \lor F) \Rightarrow F = F\);
        \item[(b)] \((T \lor F) \Rightarrow T = T\);
        \item[(c)] \((T \land F) \Rightarrow F = T\);
        \item[(d)] \((T \land F) \Rightarrow F = T\);
        \item[(e)] \(T \lor (F \Rightarrow F) = T \lor T = T\);
        \item[(f)] \(T \lor (F \Rightarrow T) = T \lor T = T\);
        \item[(g)] \(T \land (F \Rightarrow F) = T \land T = T\);
        \item[(h)] \(T \land (F \Rightarrow F) = T \land T = T\);
        \item[(i)] \((F \lor ((F \Rightarrow T) \Rightarrow T)) \Leftrightarrow (((T \Rightarrow F) \Rightarrow T) \land T)\). \\
        \(F \lor T \Leftrightarrow T \land T\) \\
        \(T \Leftrightarrow T = T\)
    \end{enumerate}

    \item Construct the truth tables for the following predicates:

    \begin{enumerate}
        \item[(a)] \(P(x) = x \lor \neg x\)
        \begin{center}
        \begin{tabular}{c|c}
        \(x\) & \(x \lor \neg x\) \\
        \hline
        T & T \\
        F & T \\
        \end{tabular}
        \end{center}

        \item[(b)] \(P(x) = x \Rightarrow \neg x\)
        \begin{center}
        \begin{tabular}{c|c}
        \(x\) & \(x \Rightarrow \neg x\) \\
        \hline
        T & F \\
        F & T \\
        \end{tabular}
        \end{center}

        \item[(c)] \(P(x, y) = x \Rightarrow \neg y\)
        \begin{center}
        \begin{tabular}{cc|c}
        \(x\) & \(y\) & \(x \Rightarrow \neg y\) \\
        \hline
        T & T & F \\
        T & F & T \\
        F & T & T \\
        F & F & T \\
        \end{tabular}
        \end{center}

        \item[(d)] \(P(x, y) = y \Rightarrow \neg x\)
        \begin{center}
        \begin{tabular}{cc|c}
        \(x\) & \(y\) & \(y \Rightarrow \neg x\) \\
        \hline
        T & T & F \\
        T & F & T \\
        F & T & T \\
        F & F & T \\
        \end{tabular}
        \end{center}

        \item[(e)] \(P(x, y, z) = (x \iff y) \land (y \iff z)\)
        \begin{center}
        \begin{tabular}{ccc|c}
        \(x\) & \(y\) & \(z\) & \(P(x, y, z)\) \\
        \hline
        T & F & F & F \\
        T & F & T & F \\
        T & T & F & F \\
        T & T & T & T \\
        F & F & F & T \\
        F & F & T & F \\
        F & T & F & F \\
        F & T & T & F \\
        \end{tabular}
        \end{center}

        \item[(f)] \(P(x, y, z) = (x \iff y) \land (y \iff z)\)
        \begin{center}
        \begin{tabular}{ccc|c}
        \(x\) & \(y\) & \(z\) & \(P(x, y, z)\) \\
        \hline
        T & F & F & F \\
        T & F & T & F \\
        T & T & F & F \\
        T & T & T & T \\
        F & F & F & T \\
        F & F & T & F \\
        F & T & F & F \\
        F & T & T & F \\
        \end{tabular}
        \end{center}

        \item[(g)] \(P(x, y, z) = (x \iff y) \land (y \iff z) \land (y \neq z)\)
        \begin{center}
        \begin{tabular}{ccc|c}
        \(x\) & \(y\) & \(z\) & \(P(x, y, z)\) \\
        \hline
        T & F & F & F \\
        T & F & T & F \\
        T & T & F & F \\
        T & T & T & F \\
        F & F & F & F \\
        F & F & T & F \\
        F & T & F & F \\
        F & T & T & F \\
        \end{tabular}
        \end{center}

        \item[(h)] \(P(x, y, z) = (x \neq y) \lor (y \neq z)\)
        \begin{center}
        \begin{tabular}{ccc|c}
        \(x\) & \(y\) & \(z\) & \(P(x, y, z)\) \\
        \hline
        T & F & F & T \\
        T & F & T & T \\
        T & T & F & T \\
        T & T & T & F \\
        F & F & F & F \\
        F & F & T & T \\
        F & T & F & T \\
        F & T & T & T \\
        \end{tabular}
        \end{center}
    \end{enumerate}



    \item \textbf{Theorem 1} (De Morgan's ``laws''). \textit{Given two statements \(P\) and \(Q\), the following equalities hold:}
\begin{gather*}
    \neg(P \lor Q) \equiv (\neg P) \land (\neg Q), \\
    \neg(P \land Q) \equiv (\neg P) \lor (\neg Q).
\end{gather*}

    Proof. It’s your task to prove this.

    Will use truth table for all possible values of P,T to prove the statements

    \begin{tabular}{|c|c|c|c|c|c|}
    \hline
    \(P\) & \(Q\) & \(\neg(P \lor Q)\) & \((\neg P) \land (\neg Q)\) & \(\neg(P \land Q)\) & \((\neg P) \lor (\neg Q)\) \\
    \hline
    T & T & F & F & F & F \\
    \hline
    T & F & F & F & T & T \\
    \hline
    F & T & F & F & T & T \\
    \hline
    F & F & T & T & T & T \\
    \hline
    \end{tabular}

    Thus, for P = \{T, F\} and Q = \{T, F\}, both statement hold


\item {Theorem 2} (De Morgan's ``laws'' for sets, universal form). \textit{Let \(U\) be a universe. For any sets \(A, B \subseteq U\) (where \(^c\) denotes complement relative to \(U\)),}
\begin{align}
(A \cup B)^c &= A^c \cap B^c, \\
(A \cap B)^c &= A^c \cup B^c.
\end{align}
\noindent \textbf{Proof.} Using Theorem 1.

LHS1. Suppose \(x \in (A \cup B)^c\). Then need to prove that \((A \cup B)^c \subseteq A^c \cap B^c\).

Proof:
\[
x \in (A \cup B)^c
\Rightarrow \neg(x \in A \cup B)
\Rightarrow \neg(x \in A \lor x \in B)
\overset{\text{Theorem 1}}{\Rightarrow} \neg(x \in A) \land \neg(x \in B)
\Rightarrow x \in A^c \cap B^c.
\]
Proven.

RHS1. Suppose \(x \in A^c \cap B^c\). Then need to prove that \(A^c \cap B^c \subseteq (A \cup B)^c\).

Proof:
\[
x \in A^c \cap B^c
\Rightarrow \neg(x \in A) \land \neg(x \in B)
\overset{\text{Theorem 1}}{\Rightarrow} \neg(x \in A \lor x \in B)
\Rightarrow \neg(x \in A \cup B)
\Rightarrow x \in (A \cup B)^c.
\]
Proven.

Then statement \((A \cup B)^c = A^c \cap B^c\) holds.

LHS2. Suppose \(x \in (A \cap B)^c\). Then need to prove that \((A \cap B)^c \subseteq A^c \cup B^c\).

Proof:
\[
x \in (A \cap B)^c
\Rightarrow \neg(x \in A \cap B)
\Rightarrow \neg(x \in A \land x \in B)
\overset{\text{Theorem 1}}{\Rightarrow} \neg(x \in A) \lor \neg(x \in B)
\Rightarrow x \in A^c \cup B^c.
\]
Proven.

RHS2. Suppose \(x \in A^c \cup B^c\). Then need to prove that \(A^c \cup B^c \subseteq (A \cap B)^c\).

Proof:
\[
x \in A^c \cup B^c
\Rightarrow \neg(x \in A) \lor \neg(x \in B)
\overset{\text{Theorem 1}}{\Rightarrow} \neg(x \in A \land x \in B)
\Rightarrow \neg(x \in A \cap B)
\Rightarrow x \in (A \cap B)^c.
\]
Proven.

Then statement \((A \cap B)^c = A^c \cup B^c\) also holds.


\item \textbf{Theorem 3} (De Morgan's ``laws'' for sets, general form). \textit{For any sets \(A, B, C\),}
\begin{align*}
A \setminus (B \cup C) &= (A \setminus B) \cap (A \setminus C), \\
A \setminus (B \cap C) &= (A \setminus B) \cup (A \setminus C).
\end{align*}

\textit{Proof.} Prove using Theorem 1.

LHS1.
Let $\Phi x \in A  \backslash (B \cup C)$.
Then $x \in A \land x \not\in (B \cup C)$ \Rightarrow $x \in A \land x \not\in B \land x \not\in C$ \Rightarrow $x \in (A \backslash B) \cap (A \backslash C)$. Proven

RHS1.
Let $\Phi x \in (A\backslash B) \cap (A \backslash C)$. Then $(x \in A \land x \not\in B) \land (x \in A \land x \not\in C)$ \Rightarrow $x \in A \land x \not\in B \land x \not\in C$ \Rightarrow $x \in A \land x \not\in (B \cup C)$ \Rightarrow $x \in (A \backslash (B \cup C))$. Proven

Then statement $A \setminus (B \cup C) &= (A \setminus B) \cap (A \setminus C)$ holds.

LHS2.
Let $\Phi x \in (A \backslash (B \cap C))$. Then $x \in A \land x \not\in (B \cap C)$ \Rightarrow $x \in A \land (x \not\in B \lor x \not\in C)$ \Rightarrow $(x \in A \land x \not\in B) \lor (x \in A \land x \not\in C)$ \Rightarrow $x \in ((A \backslash B) \cup (A \backslash C))$. Proven

RHS2.

Let $\Phi x \in ((A \backslash B) \cup (A \backslash C))$. Then $(x \in A \land x \not\in B) \lor (x \in A \land x \not\in C)$ \Rightarrow $x \in A \land (x \not\in B \lor x \not\in C)$ \Rightarrow $x \in A \land x \not\in (B \cap C)$ \Rightarrow $x \in (A \backslash (B \cap C))$. Proven

Then statement $A \setminus (B \cap C) &= (A \setminus B) \cup (A \setminus C)$ also holds.


\item \textbf{Definition 4.} \textit{Symmetrical difference} of two sets \(A\) and \(B\) is defined as a set of all elements that belong exactly to one of these two sets and is denoted as \(A \Delta B\).

\textbf{Your task:} Express the symmetrical difference in at least two different ways using all previously-defined operations \((\cup, \cap, \setminus)\). Also define it using abstraction method (it will look like \(A \Delta B = \{x \mid \dots \}\)).

$A \Delta B$ = $(A\backslash B) \cup (B \backslash A)$ = $(A \cup B)\backslash (A \cap B)$ = $\{x | (x \in A \land x \not\in B) \lor (x \not\in A \land x \in B)\}$


\end{enumerate}

\end{document}\end{document}